In this study, we systematically explore the structural stability of heptazine-based g-C$_3$N$_4$ from a monolayer to different stacking configurations and how the presence of dopants in these configurations modulates the structural and electronic properties.
Our results show that: 
\begin{itemize}
    \item 
    The widely assumed planar ($P\bar{6}m2$) model of g-C$_3$N$_4$ is neither energetically stable nor a reasonable approximation to the structure of ordered g-C$_3$N$_4$. Allowing out-of-plane atomic relaxation yields a lower-energy buckled structure that alleviates the in-plane N-lone-pair repulsion inherent to the planar model, whereas doping with heteroatoms is not required to induce buckling.
    \item 
    In the layered heptazine-based g-C$_3$N$_4$ structure, the stacking patterns previously demonstrated to be the lowest in energy for flat layers are significantly higher in energy than other stacking patterns when the layers are buckled.
    \item 
    The electronic structure of g-C$_3$N$_4$ is sensitive to both buckling and stacking registry. Although buckling generally widens the band gap relative to the planar monolayer, different corrugated stacking configurations yield distinct band gap values, which could potentially be an effective mechanism for tuning the electronic properties of bulk g-C$_3$N$_4$. For example, the corrugated stackings ($P1$ symmetry) uncovered in this study display varied electronic responses, with the band gaps of certain stacking configurations aligning closely with experimentally reported values. The most stable structure found in this work was the one labelled $AB^*_{3c}$ while $AB_{3c}$ showed the best match to the experimental band gap (calculated at the HSE06 level).
    \item 
    The calculated stability and electronic effect of representative metal and non-metal single atom dopants is strongly impacted by the treatment of both buckling and stacking.
\end{itemize}

We present a number of corrugation patterns that were calculated to be 
lower in energy than any of those previously reported.  Key to these
lower energy structures was the use of a $2 \times 2$ supercell 
within each plane, allowing adjacent heptazine units to buckle in opposite directions.  The hexagonal structure of g-C$_3$N$_4$ suggests this 
alternation of structure may be frustrated in the third equivalent in-plane direction.

Within these corrugated structures, P and Ni dopants exhibit markedly different electronic responses compared with the flat model, demonstrating that dopant functionality depends primarily on the host structural motif. As such, reliable property prediction and rational design of both pristine and doped g-C$_3$N$_4$ require careful consideration of accurate structural models to enable the development of high-performance g-C$_3$N$_4$ photocatalysts.
In particular, studying the effect of dopants by using DFT calculations based on either flat or buckled g-C$_3$N$_4$ monolayers is likely to be misleading if the target material is a synthesisable (i.e.\ multilayer or bulk) doped g-C$_3$N$_4$.  Such calculations yield incorrect predictions of both the propensity of dopants to enter the g-C$_3$N$_4$ lattice and the effects of such dopants on the electronic (and hence photocatalytic) properties of  the doped material.






